\chapter{\texorpdfstring{$\beta$}{贝塔}-消除}\label{ch:b1:elimination}

2-溴-2-甲基丙烷与乙醇钠的乙醇溶液在室温下反应时，
主要产物为醚；加热同一混合物，
则主要产物为烯烃 2-甲基丙烯，其中不含乙氧基。
乙氧根离子并非作为\omterm{def:b1:electronic-effects:nucleophile}{亲核试剂}进攻碳原子，
而是作为碱，夺取相邻碳原子上的质子，同时溴离子离去。
取代与消除之间的这种竞争贯穿整个有机合成。
本章介绍两种消除机理、
它们对几何构型的特殊要求、
预测所生成烯烃的规则，
以及如何使反应朝某一方向进行。

\begin{recall}
SN1 与 SN2（\cref{ch:b1:nucleophilic-substitution}）；\omterm{def:b1:stereochemistry-in-depth:representations}{纽曼投影式}、
环己烷的\omterm{def:b1:stereochemistry-in-depth:isomers}{构象}与 $Z$/$E$ 标记，
（\cref{ch:b1:stereochemistry-in-depth}）；\omterm{def:b1:electronic-effects:intermediates}{碳正离子}与碱，
（\cref{ch:b1:electronic-effects}）。
\end{recall}

\section{E2 机理}

\begin{definition}[\texorpdfstring{$\beta$}{贝塔}-消除、E2 与 E1]\label{def:b1:elimination:elimination}
称 \emph{$\beta$-消除}\index{$\beta$-消除}为从碳原子
（$\alpha$ 碳）上移去\omterm{def:b1:electronic-effects:nucleophile}{离去基团} Y，并从相邻碳原子（$\beta$ 碳）
上移去一个质子、两者之间形成双键的反应。
在 \emph{E2 机理}\index{E2 机理}中，碱夺取质子，
同时 Y 离去，在一个\omterm{def:b1:elementary-steps:elementary-step}{基元反应}中完成；在 \emph{E1 机理}\index{E1 机理}中，
Y 先离去，生成\omterm{def:b1:electronic-effects:intermediates}{碳正离子}，
再失去一个 $\beta$ 质子。
\end{definition}

\begin{proposition}[E2 的速率方程]\label{prop:b1:elimination:e2-law}
对于 E2，$v = k[\mathrm{RY}][\mathrm{B}]$，其中 B 为碱。
\end{proposition}

\begin{proof}
该过程为一个涉及\omterm{def:b1:nucleophilic-substitution:substitution}{底物}与碱的双分子\omterm{def:b1:elementary-steps:elementary-step}{基元反应}，由
\omterm{def:b1:elementary-steps:elementary-step}{基元反应}的质量作用定律（\cref{ch:b1:elementary-steps}）即得。
\end{proof}

\begin{definition}[反式与顺式共平面]\label{def:b1:elimination:periplanar}
相邻原子上的两条键称为\emph{反式共平面}\index{反式共平面}，
若二者位于同一平面内且分居两侧（二面角为 $180^\circ$）；
若二者位于同一平面内且位于同侧，则称为\emph{顺式共平面}\index{顺式共平面}，
此时二面角为 $0^\circ$。
\end{definition}

\begin{omfigure}
\begin{tikzpicture}[font=\small, >={Stealth}]
  \pic (n) at (0,0) {omnewman={90}{270}};
  \node at (n-f1) {H}; \node at (n-f2) {\ce{CH3}}; \node at (n-f3) {H};
  \node at (n-b1) {Br}; \node at (n-b2) {\ce{CH3}}; \node at (n-b3) {H};
  \node (b) at (-2.6,2.0) {乙醇根 \ce{EtO-}};
  \draw[->, thick, omThm] (b) to[bend left=20] ($(n-f1)+(-0.15,0.1)$);
  \draw[->, thick, omThm] (0.12,0.8) to[bend left=40] (0.12,0.25);
  \draw[->, thick, omThm] (0.12,-0.6) to[bend right=35] ($(n-b1)+(0.25,0.15)$);
  \node[below] at (0,-2.0) {H 与 Br 反式共平面};
  \draw[->, thick] (2.4,0) -- (3.6,0) node[midway, above] {E2};
  \node at (5.6,0) {\chemfig{C(-[:150]H_3C)(-[:210]H)=C(-[:30]H)(-[:-30]CH_3)}};
  \node at (5.6,-0.9) {\cip{E}-丁-2-烯};
\end{tikzpicture}

{\small E2 反应：2-溴丁烷沿 C3--C2 键观察（C3 在前，C2 带 Br 在后），
处于两种\omterm{def:b1:elimination:periplanar}{反式共平面}\omterm{def:b1:stereochemistry-in-depth:isomers}{构象}之一；另一\omterm{def:b1:stereochemistry-in-depth:isomers}{构象}中 C3 的另一个氢与 Br \omterm{def:b1:elimination:periplanar}{反式共平面}，
则生成 \cip{Z}-丁-2-烯。三对电子
同时移动：碱夺取质子，C--H 键电子对成为
$\pi$ 键，C--Br 键电子对随溴一同离去。
\omterm{def:b1:stereochemistry-in-depth:representations}{纽曼投影式}中相互正对的基团，
在双键的同一侧。}
\end{omfigure}

\begin{proposition}[E2 反式且立体专一]\label{prop:b1:elimination:anti}
E2 从 C--H 键与 C--Y 键\omterm{def:b1:elimination:periplanar}{反式共平面}的\omterm{def:b1:stereochemistry-in-depth:isomers}{构象}出发进行，
该反应具有立体专一性：非对映异构的\omterm{def:b1:nucleophilic-substitution:substitution}{底物}
生成不同的烯烃\omterm{def:b1:stereochemistry-in-depth:isomers}{立体异构体}，由反式排列所决定。
\end{proposition}

\begin{proof}
\omterm{def:b1:elementary-steps:energy-profile}{过渡态}中，C--H \omterm{def:b1:lewis-resonance-vsepr:lewis}{成键电子对}须进入 C--Y 键的
反键区，同时两个碳原子变为三角形；两键平行且位于同一平面时重叠最好。
反式排列为\omterm{def:b1:stereochemistry-in-depth:conformations}{交叉式构象}（能量较低），并使碱远离\omterm{def:b1:electronic-effects:nucleophile}{离去基团}；
顺式排列则为\omterm{def:b1:stereochemistry-in-depth:conformations}{重叠式构象}。
在 H 与 Y 反式排列时，每个碳原子上的另两个基团
位置随之确定：\omterm{def:b1:stereochemistry-in-depth:representations}{纽曼投影式}中处于同侧者，在双键上
为 \textit{顺式}\omterm{def:b1:stereochemistry-in-depth:isomers}{构型}。\omterm{def:b1:stereochemistry-in-depth:enantiomers}{非对映异构体}
在某一碳原子上有两个基团互换，因而得到另一种
烯烃。
\end{proof}

\begin{proposition}[环己烷上的 E2]\label{prop:b1:elimination:cyclohexane-e2}
在环己烷环上，E2 要求\omterm{def:b1:electronic-effects:nucleophile}{离去基团}与 $\beta$
氢均处于\omterm{def:b1:stereochemistry-in-depth:conformations}{直立键}（\textit{反式}双\omterm{def:b1:stereochemistry-in-depth:conformations}{直立键}）；
处于\omterm{def:b1:stereochemistry-in-depth:conformations}{平伏键}的\omterm{def:b1:electronic-effects:nucleophile}{离去基团}须待\omterm{def:b1:stereochemistry-in-depth:conformations}{环翻转}后方能反应。
\end{proposition}

\begin{proof}
在\omterm{def:b1:stereochemistry-in-depth:conformations}{椅式构象}中，相邻碳原子上的两条\omterm{def:b1:stereochemistry-in-depth:conformations}{直立键}\omterm{def:b1:elimination:periplanar}{反式共平面}（一条向上，
一条向下，二者平行于环轴）；\omterm{def:b1:stereochemistry-in-depth:conformations}{平伏键}仅与环上的 C--C 键反式，
从不与 C--H 键反式。
\end{proof}

\begin{omfigure}
\begin{tikzpicture}[font=\small]
  \pic (a) at (0,0) {omchair};
  \draw[thick, omThm] (a-c1) -- (a-a1) node[above] {Cl};
  \draw[thick] (a-c1) -- (a-e1) node[right] {H};
  \draw[thick, omDef] (a-c2) -- (a-a2) node[below] {H};
  \draw[thick, omDef] (a-c6) -- (a-a6) node[below] {H};
  \node[align=left, anchor=west] at (3.0,0.2) {C1 上的直立 Cl（向上）；\\两个相邻碳上的直立 H（向下）：\\两个反式共平面的 H};
\end{tikzpicture}

{\small \omterm{def:b1:stereochemistry-in-depth:conformations}{椅式构象}上的 \textit{trans}-双直立排布：直立氯
与两个相邻碳上的直立氢\omterm{def:b1:elimination:periplanar}{反式共平面}
（蓝色）。}
\end{omfigure}

\section{E1 机理}

\begin{proposition}[E1 的速率方程]\label{prop:b1:elimination:e1-law}
对 E1，$v = k[\mathrm{RY}]$，速率即电离步骤的速率。
\end{proposition}

\begin{proof}
与 SN1 相同（见\cref{prop:b1:nucleophilic-substitution:sn1-law}）：
\omterm{def:b1:electronic-effects:intermediates}{碳正离子}在慢步骤中生成、在快步骤中被消耗，
这里通过向溶剂或\omterm{def:b1:predominant-reaction:strong-acid}{弱碱}失去一个质子实现；\omterm{def:b1:elementary-steps:qssa}{稳态近似}给出
$v = k_1[\mathrm{RY}]$。
\end{proof}

E1 与 SN1 共享\omterm{def:b1:electronic-effects:intermediates}{碳正离子}：\omterm{def:b1:electronic-effects:carbon-class}{叔碳}\omterm{def:b1:nucleophilic-substitution:substitution}{底物}在\omterm{def:b1:intermolecular-forces-solvents:solvent-classes}{质子性
溶剂}中同时生成两类产物，加热有利于烯烃。由于
\omterm{def:b1:electronic-effects:intermediates}{碳正离子}为平面结构，质子随后才失去，E1
不存在几何要求，也没有立体专一性；它主要给出
更稳定的烯烃，通常为 \cip{E} 型。

\begin{omfigure}
\setchemfig{atom sep=2.1em}
\schemestart
\chemfig{C(-[2]CH_3)(-[4]H_3C)(-[6]CH_3)-Br}
\arrow{->[慢]}
\chemfig{C^{+}(-[2]CH_3)(-[:210]H_3C)(-[:-30]CH_3)}
\arrow{->[快][\ce{-H+}]}
\chemfig{C(-[2]CH_3)(-[:210]H_3C)=[:-30]CH_2}
\schemestop
\par\vspace{4pt}

{\small E1：2-溴-2-甲基丙烷电离生成\omterm{def:b1:electronic-effects:intermediates}{碳正离子}，然后
从甲基上向溶剂失去一个质子，生成
2-甲基丙烯。}
\end{omfigure}

\section{区域选择性：札依采夫规则}

\begin{definition}[区域选择性反应与立体选择性反应]\label{def:b1:elimination:selectivity}
若一个反应可生成多种\omterm{def:b1:stereochemistry-in-depth:isomers}{构造异构体}，但主要生成其中一种，
则称为\emph{区域选择性反应}\index{区域选择性反应}；
若一个反应可生成多种\omterm{def:b1:stereochemistry-in-depth:isomers}{立体异构体}，但主要生成其中一种，
则称为\emph{立体选择性反应}\index{立体选择性反应}。
\end{definition}

\begin{proposition}[札依采夫规则]\label{prop:b1:elimination:zaitsev}
当可除去多个 $\beta$ 氢时，主要烯烃通常是
取代程度最高（即最稳定）者，且为 \cip{E} 异构体
而非 \cip{Z} 异构体（\emph{札依采夫规则}\index{札依采夫规则}）。
体积庞大的碱，以及如\textit{trans}-双直立要求之类的几何约束，
可以使这一规则失效。
\end{proposition}

\begin{proof}
双键碳上的烷基使双键稳定（与它们稳定
\omterm{def:b1:electronic-effects:intermediates}{碳正离子}的作用相同），而 \cip{E} 烯烃避免了
\textit{cis} 基团的拥挤。E2 与 E1 的\omterm{def:b1:elementary-steps:energy-profile}{过渡态}
具有部分双键性质，因此更稳定的烯烃生成更快。体积庞大的
碱更易接近位阻较小的氢；而一个
不能与\omterm{def:b1:electronic-effects:nucleophile}{离去基团}\omterm{def:b1:elimination:periplanar}{反式共平面}的氢不会
被 E2 除去，无论它将生成的烯烃多么稳定。
\end{proof}

\begin{method}[预测烯烃产物]\label{met:b1:elimination:predict-alkene}
\begin{enumerate}
  \item 列出连有氢的 $\beta$ 碳。
  \item 对于 E2，只保留能与 Y \omterm{def:b1:elimination:periplanar}{反式共平面}的氢（在环上：
        即 \textit{trans}-双直立，且\omterm{def:b1:stereochemistry-in-depth:conformations}{椅式构象}确实可以形成）。
  \item 在其中优先选择生成取代程度最高烯烃的那个
        （\omterm{prop:b1:elimination:zaitsev}{札依采夫规则}），除非碱体积庞大。
  \item 对每一种情况，画出\omterm{def:b1:stereochemistry-in-depth:conformations}{对位交叉构象}的\omterm{def:b1:stereochemistry-in-depth:representations}{纽曼投影式}，
        以读出 \cip{E}/\cip{Z} \omterm{def:b1:stereochemistry-in-depth:isomers}{构型}。
\end{enumerate}
\end{method}

\section{取代还是消除？}

\begin{proposition}[取代与消除的竞争]\label{prop:b1:elimination:sn-vs-e}
强而体积庞大的碱、高温以及
\omterm{def:b1:electronic-effects:carbon-class}{叔碳}或拥挤\omterm{def:b1:nucleophilic-substitution:substitution}{底物}有利于消除；取代反应则有利于
作为\omterm{def:b1:predominant-reaction:strong-acid}{弱碱}的良好\omterm{def:b1:electronic-effects:nucleophile}{亲核试剂}、低温以及\omterm{def:b1:electronic-effects:carbon-class}{伯碳}\omterm{def:b1:nucleophilic-substitution:substitution}{底物}。
\end{proposition}

\begin{proof}
碱性强有利于进攻氢；体积大则阻碍对
拥挤碳的进攻，却不妨碍对暴露氢的进攻。\omterm{def:b1:electronic-effects:carbon-class}{叔碳}
无法被 SN2 进攻，其\omterm{def:b1:electronic-effects:intermediates}{碳正离子}则容易失去质子。消除
使一个分子变为两个或三个（烯烃、\omterm{def:b1:electronic-effects:nucleophile}{离去基团}、
质子化的碱）：温度升高时有利于消除，
这一结论在本科二年级分册中予以论证。
\end{proof}

\begin{omfigure}
\begin{tabular}{llll}
\hline
\omterm{def:b1:nucleophilic-substitution:substitution}{底物} & \omterm{def:b1:predominant-reaction:strong-acid}{强碱}，体积小 & \omterm{def:b1:predominant-reaction:strong-acid}{强碱}，体积大 & \omterm{def:b1:predominant-reaction:strong-acid}{弱碱}，良好\omterm{def:b1:electronic-effects:nucleophile}{亲核试剂} \\
\hline
\omterm{def:b1:electronic-effects:carbon-class}{伯碳} & SN2（少量 E2） & E2 & SN2 \\
\omterm{def:b1:electronic-effects:carbon-class}{仲碳} & E2（及 SN2） & E2 & SN2（非质子性），SN1/E1（质子性） \\
\omterm{def:b1:electronic-effects:carbon-class}{叔碳} & E2 & E2 & SN1/E1（质子性） \\
\hline
\end{tabular}

\smallskip
{\small 初步判断：各情形下的主要反应。加热
使每一列都更有利于消除。}
\end{omfigure}

\section{练习}

\begin{exercise}[$\star$]\label{exo:b1:elimination:1}
写出分别由 2-溴丁烷、
2-溴-2-甲基丁烷和溴代环己烷消除 HBr 可能生成的烯烃，
并指出\omterm{prop:b1:elimination:zaitsev}{札依采夫规则}预测为主要产物的那一种。
\end{exercise}

\begin{exercise}[$\star$]\label{exo:b1:elimination:2}
写出 E1 与 E2 的\omterm{def:b1:rate-laws:order}{速率方程}。对于在含乙醇钠的乙醇中的 2-溴-2-甲基丙烷，
如何通过实验区分这两种反应？
\end{exercise}

\begin{exercise}[$\star$]\label{exo:b1:elimination:3}
画出 2-溴丁烷沿 C2--C3 键的三种
\omterm{def:b1:stereochemistry-in-depth:conformations}{交叉式构象}的\omterm{def:b1:stereochemistry-in-depth:representations}{纽曼投影式}，并说明哪一种可发生 E2 反应。
\end{exercise}

\begin{exercise}[$\star$]\label{exo:b1:elimination:4}
预测下列各反应的主要反应类型（SN1、SN2、E1、E2）：溴乙烷与
氢氧化钠在水中于 \qty{25}{\celsius} 下反应；2-溴-2-甲基丙烷与
叔丁醇钾反应；2-溴丙烷与碘化钠在丙酮中反应；
2-溴-2-甲基丙烷在热乙醇中反应。
\end{exercise}

\begin{exercise}[$\star\star$]\label{exo:b1:elimination:5}
用\omterm{def:b1:stereochemistry-in-depth:representations}{纽曼投影式}证明，2-溴-3-甲基戊烷的两种\omterm{def:b1:stereochemistry-in-depth:enantiomers}{非对映异构体}中，
能失去 C3 上氢的那两种经 E2 消除 HBr，
得到 3-甲基戊-2-烯的两种\omterm{def:b1:stereochemistry-in-depth:isomers}{立体异构体}。
\end{exercise}

\begin{exercise}[$\star\star$]\label{exo:b1:elimination:6}
写出 2-甲基丁-2-醇在热的浓硫酸中脱水
的\omterm{def:b1:elementary-steps:elementary-step}{反应机理}（E1），并给出主要烯烃。
\end{exercise}

\begin{exercise}[$\star\star$]\label{exo:b1:elimination:7}
解释为什么 \textit{trans}-1-溴-4-叔丁基环己烷
以 E2 方式反应极为缓慢，而其 \textit{cis} 异构体反应很快。（叔丁基
始终保持平伏。）
\end{exercise}

\begin{exercise}[$\star\star$]\label{exo:b1:elimination:8}
2-溴-2-甲基丁烷与乙醇钠反应，主要得到
2-甲基丁-2-烯；与叔丁醇钾反应时，主要得到
2-甲基丁-1-烯。试解释原因。
\end{exercise}

\begin{exercise}[$\star\star$]\label{exo:b1:elimination:9}
为什么溴甲烷从未发生消除反应？1-溴-2,2-二甲基丙烷
经 E2 反应时又如何？
\end{exercise}

\begin{exercise}[$\star\star\star$]\label{exo:b1:elimination:10}
在 \qty{25}{\celsius} 的乙醇中，2-溴-2-甲基丙烷
生成醚（SN1）与烯烃（E1）的混合物，其比例与\omterm{def:b1:nucleophilic-substitution:substitution}{底物}浓度无关。
试证明两种产物均来自同一\omterm{def:b1:electronic-effects:intermediates}{碳正离子}，
且它们的比值是
两个\omterm{def:b1:rate-laws:order}{速率常数}之比。
\end{exercise}

\begin{exercise}[$\star\star\star$]\label{exo:b1:elimination:11}
1,2-二溴-1,2-二苯基乙烷的一种\omterm{def:b1:stereochemistry-in-depth:enantiomers}{非对映异构体}（\omterm{def:b1:stereochemistry-in-depth:enantiomers}{内消旋化合物}）
与乙醇钠发生 E2 反应，得到 1-溴-1,2-二苯基乙烯的
单一\omterm{def:b1:stereochemistry-in-depth:isomers}{立体异构体}，试用\omterm{def:b1:stereochemistry-in-depth:representations}{纽曼投影式}确定其为哪一种。
\end{exercise}

\begin{exercise}[$\star\star\star$]\label{exo:b1:elimination:12}
利用\cref{ch:b1:stereochemistry-in-depth}中的能量，解释为什么
环己烷上的 E2 反应在\omterm{def:b1:electronic-effects:nucleophile}{离去基团}必须
于拥挤\omterm{def:b1:stereochemistry-in-depth:conformations}{椅式构象}中转为\omterm{def:b1:stereochemistry-in-depth:conformations}{直立键}时，速率可减慢数百倍。
将速率表示为\omterm{def:b1:stereochemistry-in-depth:isomers}{构象}所占分数与\omterm{def:b1:rate-laws:order}{速率常数}的乘积。
\end{exercise}

\section{问题：薄荷基氯与新薄荷基氯}

\begin{problem}[{周末问题——两种非对映异构氯化物的椅式构象，
各自可用的反式共平面氢、所生成的烯烃，
以及它们的 E2 速率之比}]
\label{pb:b1:elimination:1}
薄荷基氯和新薄荷基氯分别是薄荷醇
和新薄荷醇的氯化物（见\cref{ch:b1:stereochemistry-in-depth}）。在环上，
C1 连有 Cl，C2 连有异丙基，C5 连有甲基。在薄荷基氯中，
三个基团都可以是平伏的；在新薄荷基氯中，
当两个烷基平伏时，Cl 为\omterm{def:b1:stereochemistry-in-depth:conformations}{直立键}。两者均
在乙醇中与乙醇钠反应，属于 E2 反应。题目数据：
\qty{95}{\percent} 的新薄荷基氯处于 Cl 为\omterm{def:b1:stereochemistry-in-depth:conformations}{直立键}的\omterm{def:b1:stereochemistry-in-depth:conformations}{椅式构象}中；
对于薄荷基氯，Cl 为\omterm{def:b1:stereochemistry-in-depth:conformations}{直立键}的\omterm{def:b1:stereochemistry-in-depth:conformations}{椅式构象}（三个基团均为\omterm{def:b1:stereochemistry-in-depth:conformations}{直立键}）
所占分数为 \num{5.0e-3}；假设反应性\omterm{def:b1:stereochemistry-in-depth:isomers}{构象}中每个\omterm{def:b1:elimination:periplanar}{反式共平面}氢
都以相同\omterm{def:b1:rate-laws:order}{速率常数}反应。新薄荷基氯
给出 \qty{78}{\percent} 的三取代烯烃。

\medskip
\noindent\textbf{第一部分——\omterm{def:b1:stereochemistry-in-depth:conformations}{椅式构象}。}
\begin{enumerate}
  \item 画出薄荷基氯的优势\omterm{def:b1:stereochemistry-in-depth:conformations}{椅式构象}。
  \item 画出新薄荷基氯的优势\omterm{def:b1:stereochemistry-in-depth:conformations}{椅式构象}。
  \item 这两种氯化物是\omterm{def:b1:stereochemistry-in-depth:enantiomers}{对映异构体}还是\omterm{def:b1:stereochemistry-in-depth:enantiomers}{非对映异构体}？
  \item 陈述 E2 反应在环上的几何要求。
  \item 薄荷基氯可在哪一种\omterm{def:b1:stereochemistry-in-depth:conformations}{椅式构象}中反应？画出该\omterm{def:b1:stereochemistry-in-depth:isomers}{构象}。
  \item 为什么该\omterm{def:b1:stereochemistry-in-depth:isomers}{构象}很少见？用直立基团的能量代价说明。
\end{enumerate}

\medskip
\noindent\textbf{第二部分——\omterm{def:b1:elimination:periplanar}{反式共平面}氢。}
\begin{enumerate}[resume]
  \item 列出该氯化物的 $\beta$ 碳（C2 和 C6）及其
        氢原子。
  \item 在新薄荷基氯的反应性\omterm{def:b1:stereochemistry-in-depth:isomers}{构象}中，哪些氢与 Cl
        \omterm{def:b1:elimination:periplanar}{反式共平面}？
  \item 在薄荷基氯的反应性\omterm{def:b1:stereochemistry-in-depth:isomers}{构象}中，C2 上的氢
        是否为\omterm{def:b1:stereochemistry-in-depth:conformations}{直立键}？C6 上的氢是否为\omterm{def:b1:stereochemistry-in-depth:conformations}{直立键}？
  \item 每种氯化物各提供多少个\omterm{def:b1:elimination:periplanar}{反式共平面}氢？
  \item 画出薄荷基氯反应性\omterm{def:b1:stereochemistry-in-depth:isomers}{构象}沿 C1--C6
        方向的\omterm{def:b1:stereochemistry-in-depth:representations}{纽曼投影式}。
  \item 为什么\omterm{def:b1:elimination:periplanar}{顺式共平面}
        消除不能挽救该反应？
\end{enumerate}

\medskip
\noindent\textbf{第三部分——产物。}
\begin{enumerate}[resume]
  \item 命名脱去 C2 上氢生成的烯烃，以及
        脱去 C6 上氢生成的烯烃。
  \item \omterm{prop:b1:elimination:zaitsev}{札依采夫规则}有利于哪一种烯烃？为什么？
  \item 新薄荷基氯生成什么产物？78/22 的比例
        是否与\omterm{prop:b1:elimination:zaitsev}{札依采夫规则}一致？
  \item 薄荷基氯生成什么产物？这与
        \omterm{prop:b1:elimination:zaitsev}{札依采夫规则}是否一致？
  \item 哪个反应属于\omterm{def:b1:nucleophilic-substitution:stereospecific}{立体专一反应}、\omterm{def:b1:elimination:selectivity}{区域选择性反应}，还是两者兼有？
  \item 为什么在此处用乙醇钠时，取代反应（SN2）占次要地位？
\end{enumerate}

\medskip
\noindent\textbf{第四部分——\omterm{def:b1:rate-laws:rate}{反应速率}。}
\begin{enumerate}[resume]
  \item 将每个反应的速率写成反应性\omterm{def:b1:stereochemistry-in-depth:isomers}{构象}所占分数
        与\omterm{def:b1:elimination:periplanar}{反式共平面}氢个数的函数。
  \item 计算新薄荷基氯的相对速率。
  \item 计算薄荷基氯的相对速率。
  \item 计算速率之比 $k(\text{新薄荷基})/k(\text{薄荷基})$。
\end{enumerate}
\end{problem}
